% vim: ft=tex
% SECTION: INTRODUCTION

ACRIS (Automatically-Controlled Room Illumination System) is a project I
originally envisioned in the summer of 2010.  The goal was to design a lighting
system from LEDs capable of mixing colors, install it in my dorm room, and
design algorithms for automatically controlling the color and brightness of
those LEDs.  The project went through two design iterations; the first proved
to be infeasible because it was not sufficiently modular.  The second, which I
began work on in early 2011, was born directly from the design of the Next
House Party Lighting
System,\footnote{http://next-make.mit.edu/projects/partysystem} in which I was
part of a team of five students who designed a large-scale, completely
custom-designed and built high-power lighting system for our dorm building.  As
the lead electronics designer, my primary responsibility was designing LED
controllers capable of taking serial data as input and driving high power
red-green-blue (``RGB") LEDs, white LEDs, and ultraviolet LEDs.  Our report on
this system\footnote{http://next-make.mit.edu/projects/partysystem/report.pdf}
provides details on the design of this system as well as the problems we
encountered and our solutions for them.  Figure~\ref{fig:intro:nhpls} shows the
LED controller board used in this party lighting system.

\begin{rfigure*}
    \label{fig:intro:nhpls}
    \includegraphics[scale=1.0]{fig/photo.pdf}
    %\includegraphics[scale=0.6]{fig/fig-intro-nhpls.pdf}
    % TODO nhpls board photo
    \caption{LED controller board used in the Next House Party Lighting System}
\end{rfigure*}
    
In the spring of 2011, I redesigned this board from scratch, fixing the various
issues we encountered with the older boards.  I also designed a new,
configurable power interface.  My focus was on making the board easy to use,
but also customizable, because my ultimate goal was to design a board that
would allow designers to easily implement in their projects.  In essence, I'd
like to speed up the workflow of designing custom lighting electronics.  The
board is designed for anyone who wants to build some kind of controllable light
of their own without having to design the electronics necessary for high-power
LED control.

The board I came with has the following features.

\begin{itemize}
    \item Independent control of up to 15 channels, each capable of sinking up
    to 360mA of current.
    
    \item RS-485 serial input via an RJ-45 jack (commonly used for ethernet),
    with a second jack for forwarding the signal.

    \item Customizable microcontroller code that can be extended for additional
    functionality.  A bootloader makes reprogramming the boards over the
    network very easy (e.g. multiple boards can be programmed simultaneously).

    \item Configurable power interface --- the logic and LEDs can be powered
    easily separately or by the same source.  A LM7805 can also optionally be
    used for the logic power.

    \item No common ground between devices.
\end{itemize}

This document describes the design of the board, how to build it, and how to
use it and extend it for more functionality.
